% sec-04-investigator-brochure.tex - Investigator Brochure (final IND)
% Phase 1 PDAC IND: AI Generation. IND v1.0; repository v4.3.0.
\section{Investigator Brochure}

This section is the investigational new drug (IND) level summary of the
Investigator's Brochure (IB) for the large language model (LLM) directed eight-arm
robotic pancreaticoduodenectomy with perioperative daraxonrasib Phase 1 study. It is
provided here as a clearly labeled supporting document under 21 CFR
\S 312.23(a)(5) and \S 312.55, drawing together, in one place and at IND-summary
depth, the physical, chemical, and pharmaceutical properties of the
investigational drug daraxonrasib (RMC-6236); the available nonclinical
pharmacology, pharmacokinetics, and toxicology that underpin the conservative
160 mg first-in-human starting dose; a full description of the investigational
eight-arm collaborative surgical device (56 degrees of freedom, 640 sensor
channels, the per-arm and cumulative force caps, the layered emergency-stop chain,
and the five-vessel no-fly gate); and a structured summary of the data with
explicit guidance for the investigator. The full, standalone Investigator's
Brochure is a separate bound document submitted with this IND and incorporated by
reference; the present section is the IND-level synopsis and the pointer to that
document, consistent with the ReGARDD convention of using a distinct section to
indicate and include a supporting document. The clinical-conduct detail that the
IB summarizes is carried in the study protocol (\S 6 of this IND), the chemistry,
manufacturing, and control information (\S 7), and the pharmacology and
toxicology information (\S 8), and the investigator is directed to those sections
and to the bound IB for the controlling text \cite{phase1guidance,daraxwhipple}.

The Investigator's Brochure occupies a defined place in the architecture of an
IND. Under 21 CFR \S 312.23(a)(5) the sponsor must furnish, for a study to be
conducted by an investigator other than the sponsor, a copy of the IB; and under
21 CFR \S 312.55 the sponsor must give each participating investigator a current
brochure and must keep each investigator informed of new safety information as it
emerges. The brochure is, in the language of the International Council for
Harmonisation (ICH) E6(R3) Good Clinical Practice guideline, a compilation of the
clinical and nonclinical data on the investigational article that are relevant to
its study in human subjects, assembled to provide the investigators and others
involved in the trial with the information needed to understand the rationale for,
and to comply with, the key features of the protocol. For a combined drug-device
intervention of the kind studied here the brochure must do double duty, providing
the reference safety information for both the drug evaluated under 21 CFR part 312
and the device evaluated under 21 CFR part 812, and the present section is written
to reflect that dual structure throughout \cite{ichpaistduni,cfr312paiuni}.

\subsection*{Supporting Document Status and Scope}

The Investigator's Brochure assembles, for the clinical investigator and the
reviewing institutional review board (IRB), all reference safety and efficacy
information needed to conduct the study and to assess the expectedness of adverse
events. Because the intervention has two inseparable components, an investigational
drug evaluated under 21 CFR part 312 and an investigational device evaluated under
21 CFR part 812 in a single combined IND and investigational device exemption (IDE)
submission, the IB has two coordinated parts: the daraxonrasib drug brochure and a
Physical AI device supplement that functions as the reference safety information for
the device and the on-premises advisory large language model. An adverse event is
judged unexpected, and therefore potentially reportable on an expedited timeline,
when it is not listed at the observed specificity or severity in the daraxonrasib
IB (for drug events) or in the Physical AI supplement and System Description (for
device events), as specified in \S 8 of this IND and in the protocol assessment
section. The IB is reviewed and reissued at least annually and whenever new safety
information materially changes the risk profile, and currently enrolled
participants are re-consented when such a change occurs. The version of the IB in
effect at the time of each procedure is recorded in the case report form so that
expectedness determinations are traceable to a fixed reference text
\cite{daraxwhipple,cfr312paiuni}.

The scope of this section is deliberately bounded. It is the IND-level synopsis
of the brochure and the formal pointer to the bound document; it is not the
controlling clinical-conduct text, which is the protocol (\S 6), and it is not the
controlling chemistry or nonclinical text, which is carried in the chemistry,
manufacturing, and control information (\S 7) and the pharmacology and toxicology
information (\S 8). Where this synopsis and any of those modules might appear to
differ, the bound IB and the cited module govern, and the values restated here are
restated precisely so that no discrepancy arises. The brochure and this synopsis
share one internally consistent set of quantitative facts: the three protocol dose
levels of 160 mg, 220 mg, and 300 mg once daily; the 28-day dose-limiting toxicity
(DLT) observation window; the 0.5 ng/mL serum trough threshold that governs the
perioperative pause-and-restart advisory; and the device safety envelope with its
3 N per-arm and 18 N cumulative tip-force caps, its 10 kHz heartbeat coordination
bus, its layered emergency-stop chain, and its five-vessel no-fly gate. These
values appear once in the brochure and are carried forward unchanged into every
module of the IND \cite{daraxwhipple,onpremwhippl}.

The composition of the initial IND and IRB package into which this brochure is
bound is shown in Figure 14. The clinical protocol, the Investigator's Brochure,
the nonclinical, chemistry, and stability data, the informed consent and
recruitment materials, the investigator forms (FDA Forms 1571, 1572, and 3674),
and the safety information are assembled into one internally consistent package
whose submission starts the fixed 30-day review clock of 21 CFR \S 312.40(b).
Faster, internally consistent assembly starts that clock and all participating
sites sooner, but it cannot shorten the statutory 30-day period and it cannot
generate toxicology or stability data that do not exist; the brochure is one
indispensable component of that package and is the reference safety text against
which every adverse event is later judged \cite{phase1guidance,paisitedocpk}.

\needspace{8\baselineskip}
\begin{mermaidfig}
\node[mmgoal] (pkg)  at (5.2,0)       {Initial IND + IRB package\\(highest schedule value)};
\node[mmin]   (p)    at (-1.0,2.6)    {Clinical protocol\\(Phase 1, 3+3, n=18)};
\node[mmin]   (ib)   at (3.4,2.6)     {Investigator's\\Brochure};
\node[mmin]   (nc)   at (7.8,2.6)     {Nonclinical, CMC,\\stability};
\node[mmin]   (icf)  at (-1.0,-2.6)   {ICF + recruitment\\(Physical AI opt-out)};
\node[mmin]   (form) at (3.4,-2.6)    {Investigator forms\\1571 / 1572 / 3674};
\node[mmin]   (safe) at (7.8,-2.6)    {Safety\\information};
\node[mmproc] (clk)  at (5.2,-5.4)    {30-day FDA clock\\starts at submission};
\node[mmaccent] (acc) at (1.4,-7.8)   {Faster assembly starts\\the clock and all\\sites sooner};
\node[mmdec]  (lim)  at (9.0,-7.8)    {Cannot shorten 30 days\\or create missing\\toxicology / stability};
\draw[mmedge] (p)   -- (pkg);
\draw[mmedge] (ib)  -- (pkg);
\draw[mmedge] (nc)  to[out=-90,in=40,looseness=0.5] (clk.north east);
\draw[mmedge] (icf) -- (pkg);
\draw[mmedge] (form) to[out=-90,in=140,looseness=0.5] (clk.north west);
\draw[mmedge] (safe) to[out=-90,in=60,looseness=0.5] (clk.north east);
\draw[mmedge] (pkg) -- (clk);
\draw[mmedge] (clk) -- (acc);
\draw[mmedge] (clk) -- (lim);
\end{mermaidfig}
\figcaption{Figure 14. Composition of the initial IND and IRB package. The clinical
protocol, the Investigator's Brochure, the nonclinical, chemistry, and stability
data, the informed consent and recruitment materials, the investigator forms (1571,
1572, 3674), and the safety information are assembled into one internally consistent
package. Faster, internally consistent assembly starts the 30-day FDA clock of
21 CFR \S 312.40(b) and all participating sites sooner but cannot shorten the fixed
period or generate missing toxicology or stability data.}

\subsection*{Physical, Chemical, and Pharmaceutical Properties of Daraxonrasib}

Daraxonrasib (development code RMC-6236) is an orally bioavailable, small-molecule,
RAS(ON) multi-selective (pan-RAS) inhibitor. Its pharmacological mechanism is a
tri-complex mode of action: the molecule binds the active, guanosine triphosphate
(GTP) bound conformation of mutant RAS in a ternary complex with the abundant
intracellular chaperone cyclophilin A, and in doing so it sterically and
allosterically blocks the engagement of downstream effectors, suppressing the
mitogen-activated protein kinase (MAPK) and phosphoinositide 3-kinase (PI3K)
signaling that drives pancreatic ductal adenocarcinoma (PDAC). Because it targets
the shared GTP-bound state rather than a single codon substitution, daraxonrasib is
active across the spectrum of common oncogenic KRAS variants relevant to PDAC,
including the G12D, G12V, and G12R substitutions that define the eligible population
of this study. The agent is administered as an oral tablet, once daily, and is
supplied in the strengths required to compose the three protocol dose levels of
160 mg, 220 mg, and 300 mg once daily. Tablets are swallowed whole with water and
are not crushed, split, or chewed; they are packaged in child-resistant,
light-protective containers labeled under 21 CFR \S 312.6 with the
investigational-use caution statement, the protocol number, the lot number, and the
storage conditions, and are stored at controlled room temperature in a secured,
temperature-monitored, access-controlled location until dispensed. The detailed
chemistry, manufacturing, and control characterization of the drug substance and
drug product, including appearance, composition, container-closure system, and the
stability program that governs the labeled expiry, is provided in \S 7 of this IND;
the present summary restates only the properties most relevant to the investigator
at the point of care \cite{daraxwhipple,pdac060s2030}.

The properties most material to the investigator are collected in Table 5.1. The
distinguishing pharmacological feature of daraxonrasib relative to mutation-selective
KRAS inhibitors is its breadth: a single agent addresses the three KRAS G12 variants
that together account for the great majority of PDAC, removing the need to genotype
to a single allele before treatment and matching the agent to the broad pan-KRAS
eligibility of the RASolute 302 program from which the protocol dose assumptions are
drawn. The mechanistic class, the route, the schedule, the dose levels, and the
clinical-context investigational status are stated together so that the investigator
can read the drug's identity and its place in this study from a single tabulation
\cite{daraxwhipple,natlpaiplatf}.

\needspace{10\baselineskip}
\begin{table}[t]
\caption{Daraxonrasib (RMC-6236) physical, chemical, and pharmaceutical properties material to the investigator.}
\label{tab:sec04-drugprops}
\centering
\begin{tabularx}{\textwidth}{L{4.0cm} Y}
\toprule
\textbf{Property} & \textbf{Description} \\
\midrule
Development code & RMC-6236 \\
Pharmacological class & Oral RAS(ON) multi-selective (pan-RAS) small-molecule inhibitor \\
Molecular mechanism & Tri-complex inhibitor: binds the GTP-bound active conformation of mutant RAS in a ternary complex with cyclophilin A, blocking effector engagement and suppressing MAPK and PI3K signaling \\
KRAS variant coverage & Active across common oncogenic KRAS codon-12 variants, including G12D, G12V, and G12R (the eligible PDAC genotypes) \\
Dosage form & Oral tablet, swallowed whole with water; not crushed, split, or chewed \\
Route and schedule & Oral, once daily (continuous, except across the protocol perioperative pause) \\
Protocol dose levels & 160 mg, 220 mg, and 300 mg once daily (DL1, DL2, DL3) \\
Storage and packaging & Child-resistant, light-protective containers; controlled room temperature; secured, temperature-monitored, access-controlled \\
Investigational-use labeling & 21 CFR \S 312.6 caution statement, protocol number, lot number, storage conditions \\
Regulatory standing & FDA Breakthrough Therapy (June 2025) for previously treated metastatic PDAC with KRAS G12; investigational in the perioperative curative-intent surgical context here \\
\bottomrule
\end{tabularx}
\figcaption{Table 5.1. Physical, chemical, and pharmaceutical properties of
daraxonrasib restated at IND-summary depth; the controlling chemistry,
manufacturing, and control text is in \S 7 of this IND and in the bound IB.}
\end{table}

The clinical and regulatory standing of daraxonrasib supports its use as the drug
component of this first-in-human perioperative study. The agent received FDA
Breakthrough Therapy designation in June 2025 for previously treated metastatic
PDAC harboring a KRAS G12 mutation, and the dose and exposure assumptions of this
protocol are anchored to the pan-KRAS RASolute 302 program, in which 300 mg once
daily is the established recommended Phase 2 dose (RP2D). In the present study the
drug is investigational specifically in the perioperative, curative-intent
surgical context and for its use in combination with the investigational device;
the escalation deliberately begins one and two steps below the established RP2D to
preserve a conservative first-in-human margin in this new clinical setting
\cite{daraxwhipple,natlpaiplatf}.

The dose levels and their relationship to the established RP2D are worth making
explicit because the entire first-in-human safety argument rests on them. The
three dose levels are DL1 at 160 mg once daily, DL2 at 220 mg once daily, and DL3
at 300 mg once daily, with the highest level set equal to the established RASolute
302 RP2D. The intermediate level represents an increment of 60 mg, or about
37.5 percent, above the starting dose, and the top level a further increment of
80 mg, or about 36.4 percent, above the intermediate level. The starting dose of
160 mg is therefore about 53.3 percent of the established RP2D, leaving roughly a
factor of 1.9 between first exposure in this surgical setting and the dose already
characterized for safety in the metastatic program. That margin is intentional: in
the perioperative window the drug is paused before surgery and restarted only on a
fistula-grade-keyed and trough-keyed schedule, so the conservative starting dose
protects the participant during the period of greatest physiologic stress while the
3+3 escalation establishes the maximum tolerated dose (MTD) and the RP2D in this new
context \cite{daraxwhipple,oncotrialllm}.

\subsection*{Nonclinical Pharmacology, Pharmacokinetics, and Toxicology}

The pharmacology of daraxonrasib is summarized in \S 8.1.1 of this IND and rests on
the RAS(ON) tri-complex mechanism described above: selective engagement of the
GTP-bound mutant RAS in complex with cyclophilin A, with downstream suppression of
effector signaling and consequent antiproliferative and pro-apoptotic activity in
KRAS-mutant PDAC models. Drug distribution and pharmacodynamics support a
once-daily oral schedule with sustained target coverage. The pharmacokinetic
behavior most relevant to this protocol is the relationship between the daily dose,
the serum trough concentration, and the 0.5 ng/mL trough threshold that governs the
perioperative pause-and-restart advisory. A 32-iteration deterministic
pharmacokinetic sweep, performed to characterize the perioperative window,
returned a baseline serum daraxonrasib concentration of 0.45 ng/mL at the operative
timepoint across all 32 iterations, that is, below the 0.5 ng/mL trough threshold
in every iteration, consistent with the protocol-mandated pre-operative pause; the
advisory then recommended a postoperative restart on day T+7 in 29 of 32 iterations
(90.6 percent, corresponding to a Grade A fistula with a sub-threshold trough),
on day T+14 in 3 of 32 iterations (9.4 percent, corresponding to a Grade B or
delayed serum rise), and on day T+21 in 0 of 32 iterations. This sweep is
described in the study protocol (\S 6 of this IND) and is summarized in \S 8 and
\S 3.4 \cite{qspmetpancre,fdadigtwinpc,pdacdigtwinp}.

The perioperative sweep deserves a worked reading because it is the quantitative
backbone of the restart advisory. The baseline serum concentration of 0.45 ng/mL at
the operative timepoint sits 0.05 ng/mL, or 10 percent, below the 0.5 ng/mL
threshold, and it does so in all 32 of 32 iterations rather than on average; the
sub-threshold finding is therefore robust to the parameter variation the sweep
explores, which is exactly the property a perioperative pause rule needs. The
restart-day distribution of 29 iterations at T+7, 3 iterations at T+14, and 0
iterations at T+21 maps directly onto the fistula-grade logic: the common case is a
Grade A fistula with a sub-threshold trough, which the advisory resolves to the
earliest safe restart at day T+7, while the minority case of a Grade B or
delayed-rise trajectory defers the restart to day T+14, and no iteration in this
deterministic sweep reached the T+21 tier, which is reserved for the slower
trajectories and, with a Grade C fistula, maps instead to a continued hold and to
the drug stopping rule. Table 5.2 lays the sweep out so that the advisory's behavior
can be read at a glance \cite{qspmetpancre,fdadigtwinpc,pdacdigtwinp}.

\needspace{9\baselineskip}
\begin{table}[t]
\caption{32-iteration deterministic perioperative pharmacokinetic sweep and the resulting restart-day distribution.}
\label{tab:sec04-sweep}
\centering
\begin{tabularx}{\textwidth}{L{4.6cm} C{2.4cm} C{2.0cm} Y}
\toprule
\textbf{Restart tier} & \textbf{Iterations} & \textbf{Share} & \textbf{Clinical correlate} \\
\midrule
Baseline at operative timepoint & 32 of 32 & 100\% & Serum 0.45 ng/mL, below the 0.5 ng/mL threshold in every iteration; supports the pre-operative pause \\
Restart day T+7 & 29 of 32 & 90.6\% & Grade A fistula with a sub-threshold trough; earliest safe restart \\
Restart day T+14 & 3 of 32 & 9.4\% & Grade B fistula or a delayed serum rise; restart deferred one week \\
Restart day T+21 & 0 of 32 & 0\% & Reserved for slower trajectories; with a Grade C fistula maps to a continued hold and the stopping rule \\
\bottomrule
\end{tabularx}
\figcaption{Table 5.2. The deterministic perioperative pharmacokinetic sweep
underpinning the pause-and-restart advisory. The baseline serum concentration of
0.45 ng/mL falls below the 0.5 ng/mL trough threshold in all 32 iterations, and the
restart day is keyed to the pancreaticojejunostomy ISGPS fistula grade and the serum
trough; the controlling text is the protocol (\S 6) and \S 8 of this IND.}
\end{table}

The toxicology integrated summary (\S 8.1.2) states the safety margins that
underpin the conservative DL1 160 mg first-in-human starting dose. The class
toxicity profile of RAS(ON) multi-selective inhibition is dominated by
gastrointestinal effects (diarrhea, nausea) and dermatologic effects (rash), which
are managed symptomatically and through the protocol dose-interruption and
dose-reduction rules; rescue management of these class effects, including
antidiarrheals and aggressive dermatologic care, is specified in the protocol. The
160 mg starting dose sits a full step below the 220 mg intermediate level and two
steps below the 300 mg established RP2D, providing a deliberate margin for the
perioperative surgical setting in which the drug is paused across the operative
window and restarted only on a fistula-grade-keyed and trough-keyed schedule. The
full data tabulation (\S 8.1.3) tabulates the supporting studies; for this
AI-generated submission the supporting nonclinical and in silico evidence base
includes the author's quantitative-systems-pharmacology metastatic-PDAC simulation
with verification, validation, and uncertainty quantification; the artificial
intelligence digital-twin PDAC simulations; the 100,000-patient in silico Phase 3
five-arm PDAC evaluation; and the end-to-end PDAC digital-twin trial work, each
grounded in the references identified in the figure below. The pharmacology and
toxicology integrated-summary structure that organizes \S 8 is reproduced in
Figure 13 \cite{qspmetpancre,fdadigtwinpc,pdacdigtwinp,chatgpt100kp}.

The class toxicity profile is the reference against which drug-related expectedness
is judged, so the investigator should hold it clearly in mind. The dominant events
are gastrointestinal (diarrhea and nausea) and dermatologic (rash), and these are
anticipated, dose-related, and reversible with the protocol management rules; their
appearance at the listed severities is expected and does not by itself trigger
expedited reporting, whereas an event that is not listed at the observed specificity
or severity in the daraxonrasib IB is unexpected and is assessed against the 7-day
and 15-day clocks of 21 CFR \S 312.32. Management proceeds through symptomatic
care, dose interruption, and, if a DLT-defining event occurs within the 28-day
window, the dose-reduction and 3+3 cohort rules; antidiarrheals and aggressive
dermatologic care are deployed early to keep these class effects below the severity
at which they would interrupt dosing. Table 5.3 sets out the class effects, their
expected behavior, and the management pathway so that the investigator can
distinguish an expected, manageable class effect from a reportable unexpected event
\cite{daraxwhipple,ichpaistduni}.

\needspace{9\baselineskip}
\begin{table}[t]
\caption{Daraxonrasib class toxicity profile, expectedness reference, and protocol management pathway.}
\label{tab:sec04-tox}
\centering
\begin{tabularx}{\textwidth}{L{3.4cm} L{4.0cm} Y}
\toprule
\textbf{Class effect} & \textbf{Expected behavior} & \textbf{Management and reporting reference} \\
\midrule
Diarrhea & Common, dose-related, reversible; listed in the daraxonrasib IB & Early antidiarrheals; dose interruption if severe; a DLT-defining event in the 28-day window engages the 3+3 dose-reduction rules \\
Nausea & Common, dose-related; listed in the IB & Antiemetic prophylaxis and treatment; dose interruption if severe; graded by CTCAE v5.0 \\
Rash (dermatologic) & Common, dose-related; listed in the IB & Aggressive dermatologic care begun early; dose interruption if severe; reduction per protocol \\
Unlisted event & Not listed at the observed specificity or severity & Unexpected; assess against 21 CFR \S 312.32 (7-day fatal or life-threatening; 15-day other serious and unexpected) \\
\bottomrule
\end{tabularx}
\figcaption{Table 5.3. The drug class toxicity profile that serves as the
expectedness reference for daraxonrasib. Gastrointestinal and dermatologic events at
the listed severities are expected and managed per protocol; an event not listed at
the observed specificity or severity is unexpected and is assessed against the
regulatory reporting clocks.}
\end{table}

\needspace{8\baselineskip}
\begin{mermaidfig}
\node[mmgoal] (pt)    at (5.2,0)      {Pharmacology and\\Toxicology (\S 8)};
\node[mmproc] (ph)    at (0,-2.6)     {8.1.1 Pharmacology\\summary and\\conclusions};
\node[mmproc] (tx)    at (5.2,-2.6)   {8.1.2 Toxicology:\\integrated summary};
\node[mmproc] (ft)    at (10.4,-2.6)  {8.1.3 Toxicology:\\full data tabulation};
\node[mmin]   (pd)    at (-3.2,-5.4)  {Mechanism,\\distribution, RAS(ON)\\pharmacodynamics};
\node[mmin]   (dsim)  at (1.2,-5.4)   {Digital-twin and QSP\\simulation evidence};
\node[mmin]   (safem) at (5.6,-5.4)   {Safety margins; DL1\\160 mg conservative\\first-in-human};
\node[mmin]   (sweep) at (10.4,-5.4)  {32-iteration\\perioperative PK sweep\\(0.45 ng/mL baseline)};
\draw[mmedge] (pt) -- (ph);
\draw[mmedge] (pt) -- (tx);
\draw[mmedge] (pt) -- (ft);
\draw[mmedge] (ph) -- (pd);
\draw[mmedge] (ph) to[out=-35,in=120,looseness=0.6] (dsim.north);
\draw[mmedge] (tx) -- (safem);
\draw[mmedge] (ft) -- (sweep);
\end{mermaidfig}
\figcaption{Figure 13. The pharmacology and toxicology integrated-summary structure
for \S 8. The pharmacology summary covers the RAS(ON) mechanism, drug distribution,
and pharmacodynamics; the integrated toxicology summary states the safety margins
underpinning the conservative DL1 160 mg first-in-human starting dose; and the full
data tabulation tabulates the supporting studies, here grounded in the author's
digital-twin and quantitative-systems-pharmacology simulation evidence and the
32-iteration perioperative pharmacokinetic sweep (0.45 ng/mL baseline serum trough
at the operative timepoint across all iterations).}

The supporting in silico evidence base named in the full data tabulation deserves to
be enumerated because it is the principal nonclinical foundation of this
AI-generated submission. The quantitative-systems-pharmacology metastatic-PDAC
simulation supplies a mechanistic model of drug exposure and effect with explicit
verification, validation, and uncertainty quantification; the artificial-intelligence
digital-twin PDAC simulations supply patient-level trajectory modeling that informs
the perioperative window; the 100,000-patient in silico Phase 3 five-arm PDAC
evaluation supplies a large-scale comparative-effectiveness context; and the
end-to-end PDAC digital-twin trial work supplies the integrated trial-simulation
framework from which the perioperative sweep is drawn. Each is a computational rather
than a wet-laboratory study, and the brochure presents it as such; the reviewer is
directed to \S 8.1.3 for the full tabulation and to the cited records for the
underlying methods and results \cite{qspmetpancre,fdadigtwinpc,chatgpt100kp,pdacdigtwinp}.

\subsection*{Investigational Device Description: the Eight-Arm Robotic System}

The investigational device is an on-premises, LLM-directed, eight-arm parallel
coelomic surgical platform that conducts the pancreaticoduodenectomy and generates
the operative and pathology inputs to the perioperative drug-restart advisory. The
on-premises, repository-pinned LLM acts strictly as a second-opinion advisory
oracle held outside the robot vendor kinematic stack and is never an autonomous
controller; full autonomy is prohibited consistent with 21 CFR \S 312.21(e), and
in this Phase 1 study the device operates only at the clinician-controlled
teleoperated and supervised semiautonomous autonomy levels under continuous human
oversight with immediate manual override. The platform carries 56 mechanical
degrees of freedom, composed of eight arms with seven degrees of freedom each
(8 multiplied by 7), and a 640-channel mixed-rate sensor stack of 80 channels per
arm. Force channels sample at 100 kHz and all remaining channels sample at the
10 kHz command rate. Positioning accuracy is 0.05 mm root mean square at a peak tip
velocity of 1,200 mm/s. The cyber-physical safety envelope comprises a per-arm
tip-force cap of no more than 3 N and a cumulative cross-arm tip-force cap of no
more than 18 N enforced at the heartbeat boundary; a 10 kHz heartbeat coordination
bus broadcasting a 64-byte frame on a 100 microsecond per-arm watchdog deadline; an
emergency arm park within 50 microseconds of a missed or out-of-parameter frame; a
cross-arm emergency stop (E-stop) that halts the full platform within 3 ms at a
0.0 N force clamp; and a software-independent hardware E-stop that meets the 21 CFR
\S 312.404 requirement of no more than 500 ms. The device specifications are
summarized in Table 5.4, and the layered safety envelope, including the five-vessel
no-fly gate, is summarized in Table 5.5 \cite{daraxwhipple,onpremwhippl}.

The numeric structure of the platform is internally consistent and worth stating as
an arithmetic identity: the 56 mechanical degrees of freedom are exactly eight arms
multiplied by seven degrees of freedom per arm, and the 640 sensor channels are
exactly eight arms multiplied by 80 channels per arm. The mixed-rate sensing
partitions cleanly: the force channels, which carry the per-arm and cumulative
tip-force signals that the safety caps police, sample at 100 kHz, a factor of ten
faster than the 10 kHz command rate at which the remaining channels and the
heartbeat coordination bus operate. That ten-to-one ratio is deliberate, because
force is the physical quantity whose excursion must be caught fastest. The
positioning accuracy of 0.05 mm root mean square is maintained at a peak tip velocity
of 1,200 mm/s, and both are verified in the pre-procedure safety matrix at the start
of each case rather than assumed from a prior calibration \cite{daraxwhipple,onpremwhippl}.

\needspace{12\baselineskip}
\begin{table}[t]
\caption{Eight-arm LLM-directed robotic Whipple device specifications (Investigator's Brochure device supplement; 21 CFR \S 312.23(g) hardware fields).}
\label{tab:sec04-device}
\centering
\begin{tabularx}{\textwidth}{L{4.2cm} L{4.4cm} Y}
\toprule
\textbf{Specification} & \textbf{Value} & \textbf{Notes} \\
\midrule
Architecture & Eight-arm parallel coelomic platform & On-premises LLM advisory oracle outside vendor kinematic stack \\
Degrees of freedom & 56 total (8 arms x 7 DOF) & Cooperating arms, phase-keyed tool assignment P1 to P8 \\
Sensor channels & 640 total (80 per arm) & Force at 100 kHz; all other channels at 10 kHz command rate \\
Positioning accuracy & 0.05 mm RMS at 1,200 mm/s & Verified in the pre-procedure safety matrix each case \\
Per-arm tip-force cap & 3 N maximum & Enforced continuously per arm \\
Cumulative force cap & 18 N maximum & Enforced at the heartbeat boundary across all arms \\
Heartbeat bus & 10 kHz, 64-byte frame & 100 microsecond per-arm watchdog deadline \\
Emergency arm park & 50 microseconds maximum & On a missed or out-of-parameter frame \\
Cross-arm E-stop & 3 ms maximum, 0.0 N clamp & Software escalation halting the full platform \\
Hardware E-stop & 500 ms maximum & Software-independent; meets 21 CFR \S 312.404 \\
Autonomy level (Phase 1) & Teleoperated / supervised & Continuous human oversight; full autonomy prohibited (\S 312.21(e)) \\
\bottomrule
\end{tabularx}
\figcaption{Table 5.4. Device specifications carried in the Investigator's Brochure
device supplement and the Physical AI System Description; the controlling
hardware, software, and safety-system text is in \S 6 and \S 7 of this IND and in
the bound IB.}
\end{table}

The eight cooperating arms are coordinated by the 10 kHz heartbeat bus and carry
phase-keyed tool assignments across the eight operative phases. Phases P1 through
P4 comprise the dissection and specimen removal: Arms 1 and 2 perform the active
dissection under the vessel control of Arm 3 (bipolar plus retractor) and the
near-infrared indocyanine-green imaging of Arm 4, with Arm 7 providing suction and
irrigation throughout. The specimen is removed at the end of P4 and reconstruction
begins. The three reconstructive anastomoses are then constructed in sequence under
closed-loop ring-tension control by the Arm 5 anastomosis probe, which issues a
soft warning whenever measured ring tension falls outside its controller band: the
duct-to-mucosa pancreaticojejunostomy (PJ) in P5 with a band of 0.30 to 0.60 N, the
end-to-side hepaticojejunostomy (HJ) in P6 with a band of 0.20 to 0.50 N, and the
antecolic gastrojejunostomy (GJ) in P7 with a band of 0.40 to 0.80 N; P8 covers
final imaging, hemostasis, and drain placement. The pancreaticojejunostomy is the
dominant determinant of postoperative pancreatic-fistula risk and is therefore the
anastomosis whose ISGPS fistula grade supplies the input to the perioperative
drug-restart advisory \cite{daraxwhipple,paisitedocpk}.

The phase-keyed tool assignment and the per-anastomosis ring-tension bands are
collected in Table 5.6 so that the investigator can see the operative choreography
and its closed-loop controls in one place. The three ring-tension bands are
anastomosis-specific because the three reconstructions differ in tissue and in
consequence: the pancreaticojejunostomy band of 0.30 to 0.60 N reflects the soft,
friable pancreatic remnant and the fistula risk that makes this anastomosis the
keyed input to the restart advisory; the hepaticojejunostomy band of 0.20 to 0.50 N
reflects the thinner-walled biliary tissue; and the gastrojejunostomy band of 0.40
to 0.80 N reflects the more robust gastric and jejunal walls. A measured tension
outside the controller band raises a soft warning from the Arm 5 anastomosis probe
rather than halting the platform, because a tension excursion is an advisory signal
to the operating surgeon, not a vascular-proximity hazard; the hard halts are
reserved for the no-fly gate and the heartbeat watchdog described below
\cite{daraxwhipple,onpremwhippl}.

\needspace{10\baselineskip}
\begin{table}[t]
\caption{Phase-keyed tool assignment and closed-loop ring-tension control across the eight operative phases.}
\label{tab:sec04-phases}
\centering
\begin{tabularx}{\textwidth}{L{2.0cm} L{5.4cm} Y}
\toprule
\textbf{Phase} & \textbf{Activity and arm assignment} & \textbf{Closed-loop control} \\
\midrule
P1 to P4 & Dissection and specimen removal: Arms 1 and 2 dissect under Arm 3 vessel control (bipolar plus retractor) and Arm 4 near-infrared ICG imaging; Arm 7 suction and irrigation throughout & Five-vessel no-fly gate active at 10 kHz; force caps enforced; specimen removed at the end of P4 \\
P5 (PJ) & Duct-to-mucosa pancreaticojejunostomy by the Arm 5 anastomosis probe & Ring-tension band 0.30 to 0.60 N; soft warning outside band; ISGPS grade keys the restart advisory \\
P6 (HJ) & End-to-side hepaticojejunostomy & Ring-tension band 0.20 to 0.50 N; soft warning outside band \\
P7 (GJ) & Antecolic gastrojejunostomy & Ring-tension band 0.40 to 0.80 N; soft warning outside band \\
P8 & Final imaging, hemostasis, and drain placement & Imaging confirmation; no-fly gate and force caps remain active \\
\bottomrule
\end{tabularx}
\figcaption{Table 5.6. The eight operative phases, the phase-keyed arm and tool
assignment, and the per-anastomosis closed-loop ring-tension bands. The
pancreaticojejunostomy (P5) is the dominant fistula-risk determinant and supplies
the ISGPS grade to the perioperative drug-restart advisory.}
\end{table}

\subsection*{Safety Envelope: Force Caps, E-Stop Chain, and the No-Fly Gate}

The device safety envelope is enforced before and during every procedure rather
than reconstructed afterward. The vascular safety-zone gate ticks at 10 kHz and
evaluates the instrument-tip proximity to each of five named vessels, the superior
mesenteric vein (SMV), the portal vein (PV), the hepatic artery (HA), the celiac
axis (CA), and the superior mesenteric artery (SMA), against three concentric
thresholds: a soft-warning zone at 3.0 mm that raises an advisory but does not block
the command, a no-fly zone that blocks the command (1.0 mm at the SMV and PV; 1.5 mm
at the HA, CA, and SMA), and a hard-stop zone at 5.0 mm that forces a cross-arm
E-stop within 3 ms. Across a representative 100-tick verdict sample the gate
returned clear on 83 ticks, soft warning on 6, no-fly on 6, and hard-stop on 5,
exercising all five vessels across all eight operative phases. The halt chain that
backs the gate is layered: an on-time heartbeat frame continues the command state,
whereas a missed or out-of-parameter frame parks the affected arm within 50
microseconds and escalates to the cross-arm E-stop, and the software-independent
hardware E-stop backs the software chain at the 500 ms outer bound. The safety
envelope is summarized in Table 5.5 \cite{daraxwhipple,clintrustpai}.

The geometry of the no-fly gate is worth stating precisely because it is
asymmetric by design. The soft-warning radius of 3.0 mm and the hard-stop radius of
5.0 mm are uniform across all five vessels, but the blocking no-fly radius is
vessel-specific: 1.0 mm at the two veins (SMV and PV) and the larger 1.5 mm at the
three arteries (HA, CA, and SMA). The arterial margin is set wider because an
arterial injury is the more catastrophic event, so the gate blocks the command at a
greater standoff for the hepatic artery, the celiac axis, and the superior
mesenteric artery than for the two veins. The ordering of the three radii is fixed
and monotone, with the hard-stop radius (5.0 mm) outside the soft-warning radius
(3.0 mm), which in turn lies outside the blocking radius (1.0 or 1.5 mm); a tip
approaching a vessel therefore crosses the hard-stop boundary first and triggers an
E-stop before it could reach the blocking or warning radii, which is the
conservative ordering a vascular-safety gate requires. The representative 100-tick
verdict sample partitions exactly into 83 clear, 6 soft warning, 6 no-fly, and 5
hard-stop ticks, summing to 100, and exercises all five vessels across all eight
operative phases \cite{daraxwhipple,clintrustpai}.

\needspace{11\baselineskip}
\begin{table}[t]
\caption{Device safety envelope: force caps, layered emergency-stop chain, and the five-vessel no-fly gate.}
\label{tab:sec04-safety}
\centering
\begin{tabularx}{\textwidth}{L{3.6cm} L{4.6cm} Y}
\toprule
\textbf{Safety element} & \textbf{Threshold / action} & \textbf{Basis and verdict sample} \\
\midrule
Per-arm force cap & 3 N maximum per arm & Continuous; force channels at 100 kHz \\
Cumulative force cap & 18 N maximum across arms & Enforced at the heartbeat boundary \\
Heartbeat and watchdog & 10 kHz bus, 100 microsecond deadline & On-time frame continues command state \\
Arm park & 50 microseconds maximum & Triggered by a missed or out-of-parameter frame \\
Cross-arm E-stop & 3 ms maximum, 0.0 N clamp & Software escalation; halts full platform \\
Hardware E-stop & 500 ms maximum & Independent of software; 21 CFR \S 312.404 \\
Soft-warning zone & 3.0 mm, all five vessels & Advisory only; 6 of 100 ticks \\
No-fly zone (block) & 1.0 mm SMV/PV; 1.5 mm HA/CA/SMA & Command blocked; 6 of 100 ticks \\
Hard-stop zone & 5.0 mm, all five vessels & Forces E-stop; 5 of 100 ticks \\
Clear verdict & Outside all zones & Command proceeds; 83 of 100 ticks \\
\bottomrule
\end{tabularx}
\figcaption{Table 5.5. The cyber-physical safety envelope and the five-vessel
no-fly gate (SMV, PV, hepatic artery, celiac axis, SMA), sampled at 10 kHz; the
representative 100-tick verdict sample partitions into 83 clear, 6 soft warning,
6 no-fly, and 5 hard-stop ticks. The controlling text is in the protocol assessment
section (\S 6 and \S 8 of this IND).}
\end{table}

Every advisory command, every safety-gate verdict, every force and trajectory
sample at the recorded rates, and every E-stop event is written to a 21 CFR part 11
hash-chained, append-only audit trail, so that compliance with the cyber-physical
envelope is verifiable for each case and any deviation is time-stamped and
detectable. The audit trail also satisfies the preservation requirement from 24
hours before to 72 hours after any reportable Physical AI adverse event. Before any
patient-facing robotic activity, the platform must clear the Phase 0 simulation
validation gate (a Unified Safety Level of at least 7.0; at least 1,000 simulations
across at least two simulation frameworks; a sim-to-real trajectory gap below 2 mm;
and a tip-force gap below 0.5 N), and the gate is re-cleared after any change. The
on-premises LLM is hash-pinned to a deterministic seed and repository commit and is
governed by a verification-before-generation gate suite, aligned with H.R. 9510,
the Verification Before Generation in Physical AI Oncology Trials Act of 2026
\cite{congress9510,clintrustpai,paiautospons}.

The Phase 0 simulation validation gate is a hard precondition rather than a target,
and its four criteria are conjunctive: the Unified Safety Level must reach at least
7.0; at least 1,000 simulations must be run across at least two independent
simulation frameworks; the sim-to-real trajectory gap must fall below 2 mm; and the
sim-to-real tip-force gap must fall below 0.5 N. All four must hold simultaneously
before any patient-facing robotic activity, and the gate is re-cleared in full after
any change to the platform, the model, or the safety configuration, so that no
modification is allowed to reach a patient without a fresh validation. The
sim-to-real tolerances connect directly to the Physical AI reporting trigger for
sim-to-real divergence: a trajectory gap beyond twice the 2 mm tolerance (that is,
above 4 mm) or a tip-force gap beyond twice the 0.5 N tolerance (that is, above
1.0 N) is a reportable divergence under the six Physical AI triggers of 21 CFR
\S 312.32(g). The investigator should read the Phase 0 gate, the run-time safety
envelope, and the reporting triggers as one continuous chain of defenses, each
checked against a fixed numeric criterion \cite{congress9510,clintrustpai}.

\subsection*{Summary of Data and Guidance for the Investigator}

The investigator should treat the daraxonrasib drug brochure and the Physical AI
device supplement as the two controlling reference texts for expectedness and risk.
For the drug, the established 300 mg once-daily RP2D and the class toxicity profile
(gastrointestinal and dermatologic effects) define the expected events; the study
deliberately starts at 160 mg once daily to preserve a conservative first-in-human
margin in the perioperative setting, and the 28-day dose-limiting toxicity window,
the 3+3 escalation, and the dose-interruption and dose-reduction rules govern
management. For the device, the force caps, the layered E-stop chain, and the
five-vessel no-fly gate define the expected safety behavior; any robotic
malfunction, AI prediction error, sensor failure, communication loss, unauthorized
access, digital-twin divergence, or event requiring activation of the E-stop is a
Physical AI adverse event under 21 CFR \S 312.32(f) and is assessed against this
supplement for expectedness. The perioperative pause-and-restart advisory is
LLM-bound and human-confirmed and never autonomous; the investigator confirms the
restart day keyed to the pancreaticojejunostomy ISGPS fistula grade and the serum
trough relative to 0.5 ng/mL, with a Grade C fistula mapping to T+21 or a continued
hold and to the drug stopping rule. All adverse events are graded by CTCAE version
5.0 and reported on the regulatory timelines (no later than 7 calendar days for an
unexpected fatal or life-threatening suspected adverse reaction and no later than
15 calendar days for any other serious and unexpected suspected adverse reaction),
with the six Physical AI reporting triggers of \S 312.32(g) layered on top
\cite{daraxwhipple,cfr312paiuni,ichpaistduni}.

The six Physical AI reporting triggers of 21 CFR \S 312.32(g) layered on top of the
ordinary drug-safety clocks are stated in Table 5.7 so that the investigator can
match an observed device or model event to the correct reporting clock without
ambiguity. The first trigger is a serious Physical AI adverse event, reported on the
ordinary 7-day or 15-day clock by seriousness and expectedness; the second is a
system-drug interaction, reported within 15 days; the third is a cybersecurity
incident, reported within 15 days or within 7 days if there is potential for patient
harm; the fourth is a sustained AI model degradation, defined as degradation across
at least three consecutive procedures or at least 24 cumulative hours; the fifth is a
sim-to-real divergence beyond twice the tolerance (trajectory above 4 mm or tip-force
above 1.0 N); and the sixth is a digital-twin discrepancy. These triggers operate in
addition to, not instead of, the ordinary serious-adverse-reaction reporting of
\S 312.32(c), and the device supplement is the expectedness reference for each
\cite{cfr312paiuni,clintrustpai,paiautospons}.

\needspace{11\baselineskip}
\begin{table}[t]
\caption{The six Physical AI reporting triggers of 21 CFR \S 312.32(g) and their clocks, layered on the ordinary drug-safety reporting.}
\label{tab:sec04-triggers}
\centering
\begin{tabularx}{\textwidth}{L{0.7cm} L{4.6cm} Y}
\toprule
\textbf{\#} & \textbf{Trigger} & \textbf{Reporting clock and definition} \\
\midrule
1 & Serious Physical AI adverse event & 7 days (unexpected fatal or life-threatening) or 15 days (other serious and unexpected), as for a drug suspected adverse reaction \\
2 & System-drug interaction & Within 15 days \\
3 & Cybersecurity incident & Within 15 days, or within 7 days if there is potential for patient harm \\
4 & Sustained AI model degradation & Reportable; degradation across at least 3 consecutive procedures or at least 24 cumulative hours \\
5 & Sim-to-real divergence & Divergence beyond twice the tolerance: trajectory above 4 mm or tip-force above 1.0 N \\
6 & Digital-twin discrepancy & Reportable discrepancy between the digital twin and the observed system \\
\bottomrule
\end{tabularx}
\figcaption{Table 5.7. The six Physical AI reporting triggers of 21 CFR
\S 312.32(g), layered on top of the ordinary 7-day and 15-day drug-safety clocks.
The Physical AI device supplement of the Investigator's Brochure is the
expectedness reference against which each trigger is judged; the audit trail is
preserved from 24 hours before to 72 hours after any reportable Physical AI adverse
event.}
\end{table}

The categories of event that constitute a Physical AI adverse event under 21 CFR
\S 312.32(f) are enumerated so that the investigator does not have to infer them.
A robotic malfunction, an AI prediction error, a sensor failure, a communication
loss, an instance of unauthorized access, a digital-twin divergence, and any event
that requires activation of the emergency stop each constitute a Physical AI adverse
event and each is assessed against the device supplement for expectedness. The
audit trail is the evidentiary record for that assessment: because every advisory
command, every gate verdict, every force and trajectory sample, and every E-stop
event is hash-chained and append-only, the sequence of events around any reportable
occurrence can be reconstructed exactly, and the preservation window from 24 hours
before to 72 hours after the event ensures the relevant record survives. The
three-tier oversight structure (the data safety monitoring board with halt
authority, the independent safety monitor with per-procedure real-time stop
authority, and the Physical AI Safety Review Committee on a cadence of no more than
90 days) reviews these events, and the binding halt rules (one or more device or
procedure serious adverse events within a cohort, or two device-related deaths at
any time) trigger a board review and possible halt \cite{cfr312paiuni,clintrustpai,paiautospons}.

The investigator is reminded that the controlling clinical-conduct text is the
study protocol and that the present section is an IND-level summary; the full
standalone Investigator's Brochure is a separate bound document submitted with this
IND and is the authoritative reference safety information for both the drug and the
device. The investigator should hold a current copy of the bound IB, record its
version in the case report form for each procedure, and consult \S 6, \S 7, and
\S 8 of this IND for the drug-and-device intervention detail, the chemistry,
manufacturing, and control information, and the pharmacology and toxicology
information that the IB summarizes. The brochure is reviewed and reissued at least
annually and whenever new safety information materially changes the risk profile,
and currently enrolled participants are re-consented when such a change occurs, so
that the reference text against which every expectedness determination is made
remains current throughout the conduct of the study
\cite{phase1guidance,daraxwhipple,oncotrialllm}.
